\section{Prompts}\label{app:prompts}

Every model call consists of a system prompt and one user message.
The system prompt is fixed for each condition; the user message is assembled from the current state at every step.
When a reply cannot be read, one retry message is appended and the model is called again; a second failure ends the step.
All texts below were captured from the requests the harness sends, not copied from the source code.

\subsection{Agent episodes}\label{app:prompts:episodes}

J, J+ and H share the selection prompt: the model answers with the number of an entry in a list.
\begin{Verbatim}
You are a web agent driving a website by choosing among the actions it currently offers.
You are given a goal and a numbered list of AVAILABLE ACTIONS. Each action has a method, a url, and zero or more input fields.
Choose the SINGLE action that best advances the goal right now.
Provide values ONLY for fields whose value you can derive from the goal. Never invent personal data (names, phone numbers) that is not in the goal — leave those out and they will be requested from the user.
If the goal is already achieved, or no available action helps, use index -1.
Respond with ONLY a JSON object: {"index": <number>, "values": {"<fieldName>": "<value>"}}. No prose, no code fences.
\end{Verbatim}
Retry message:
\begin{Verbatim}
Return ONLY a single JSON object {"index":..,"values":{..}} and nothing else.
\end{Verbatim}
P gets the construction prompt: there is no list, and the model writes the request itself.
\begin{Verbatim}
You are a web agent driving a website. You are given a goal and the current page's HTML.
There is no list of available actions: work out for yourself what request to send next.
Read the page, decide what to do, and emit the HTTP request a browser would send.
Provide values ONLY for parameters you can derive from the goal or from the page. Never invent personal data (names, phone numbers) that is not in the goal — leave those out and they will be requested from the user.
If the goal is already achieved, or nothing on this page helps, use method "STOP".
Respond with ONLY a JSON object: {"method": "GET|POST|PUT|DELETE|STOP", "url": "/path", "values": {"<param>": "<value>"}}. No prose, no code fences.
\end{Verbatim}
Retry message:
\begin{Verbatim}
Return ONLY a single JSON object {"method":..,"url":..,"values":{..}} and nothing else.
\end{Verbatim}
X and X+ get the tool-calling prompt. The tools are passed in the \texttt{tools} parameter of the request, not in the text (\autoref{app:surfaces}).
\begin{Verbatim}
You are a web agent driving a website by calling the tools it currently offers.
You are given a goal and the current page. Call the SINGLE tool that best advances the goal right now.
Provide values ONLY for arguments you can derive from the goal. Never invent personal data (names, phone numbers) that is not in the goal — leave those out and they will be requested from the user.
If the goal is already achieved, or no available tool helps, do NOT call a tool: reply with the single word STOP.
\end{Verbatim}
Retry message:
\begin{Verbatim}
Call exactly one of the available tools, or reply with the single word STOP.
\end{Verbatim}
The user message has the same blocks in every condition.
\texttt{ACTIONS TAKEN SO FAR} appears from the second step on; P calls it \texttt{REQUESTS SENT SO FAR} and receives the raw HTML instead of the page text and the list.
\begin{Verbatim}
GOAL: <goal text of the task>

CURRENT PAGE:
<visible text of the current page>

ACTIONS TAKEN SO FAR:
- <one line per earlier step>

AVAILABLE ACTIONS:
<numbered list; J, J+ and H only>
\end{Verbatim}

\subsection{Interface audit}\label{app:prompts:audit}

In the audit the model designs the controls for one screen instead of acting.
All conditions share the first part of the system prompt:
\begin{Verbatim}
You design the interface a user sees for the current step of a task on a website.
An assistant is helping the user. You decide which buttons and forms to put on screen.
Show the controls that are relevant to the goal and that the user may want at this step, not only the single next action.
Show only controls the user can actually use on this site right now. Leave everything else out.
If nothing on the site can help with the goal right now, show no controls and say so in the message.
For each field of a control, say whether the interface fixes its value ("fixed"), fills in a value the user may change ("prefill"), or asks the user ("ask").
Never invent personal data (names, phone numbers, card numbers) that is not in the goal or on the page. Mark such fields "ask".
\end{Verbatim}
The second part says how a control refers to the site. Under H:
\begin{Verbatim}
Each control must be one of the AVAILABLE CONTROLS below, referenced by its number. Their fields are drawn automatically; do not list them.
Respond with ONLY a JSON object: {"message": "<text for the user>", "controls": [{"ref": <number>, "label": "<button text>"}]}
\end{Verbatim}
Under J and J+ (\texttt{OPERATIONS}) and X+ (\texttt{TOOLS}):
\begin{Verbatim}
Each control invokes one of the OPERATIONS below, referenced by its name.
Respond with ONLY a JSON object: {"message": "<text for the user>", "controls": [{"ref": "<name>", "label": "<button text>", "fields": [{"name": "<parameter>", "mode": "fixed|prefill|ask", "value": "<value or omit>"}]}]}
\end{Verbatim}
Under P:
\begin{Verbatim}
Each control is an HTTP request this site accepts, written as method and path.
Respond with ONLY a JSON object: {"message": "<text for the user>", "controls": [{"ref": {"method": "GET|POST|PUT|DELETE", "url": "/path"}, "label": "<button text>", "fields": [{"name": "<parameter>", "mode": "fixed|prefill|ask", "value": "<value or omit>"}]}]}
\end{Verbatim}
Retry message:
\begin{Verbatim}
Return ONLY the JSON object described, nothing else.
\end{Verbatim}
The user message has the goal and the page text as in the episodes, followed by \texttt{AVAILABLE CONTROLS} (H), \texttt{OPERATIONS} (J, J+) or \texttt{TOOLS} (X+); under P it has the raw HTML only.
In the audit the X+ tools are listed as text, because generating an interface involves no tool call.

\clearpage

\section{Task Set}\label{app:tasks}

The rules by which the tasks are derived from the state graph are described in \autoref{sec:eval:method:tasks}.

{\footnotesize
\begin{longtable}{@{}lP{0.19\textwidth}P{0.24\textwidth}P{0.2\textwidth}P{0.15\textwidth}@{}}
\caption{All 24 tasks. \emph{Start}: status, party size, date and pre-order of the start state. \emph{Satisfied by}: the operations that achieve the goal. \emph{Route}: for T2, the request that leads to a state where one of them is valid. Generated from \texttt{exp/tasks.json}.}
\label{tab:app:tasks} \\
\hline
Task & Start & Goal & Satisfied by & Route \\
\hline
\endfirsthead
\hline
Task & Start & Goal & Satisfied by & Route \\
\hline
\endhead
\hline
\endfoot
T1-01 & \texttt{deposit\_\allowbreak{}pending}, party 8, \texttt{next-week} & ``Go ahead and confirm the booking.'' & \texttt{confirm\_\allowbreak{}booking}, \texttt{pay\_\allowbreak{}deposit} & -- \\
T1-02 & \texttt{confirmed}, party 8, \texttt{tomorrow} & ``We cannot make it — cancel the booking.'' & \texttt{cancel\_\allowbreak{}booking}, \texttt{request\_\allowbreak{}cancellation}, \texttt{discard\_\allowbreak{}draft}, \texttt{release\_\allowbreak{}hold} & -- \\
T1-03 & \texttt{browsing} & ``Find out which tables are free.'' & \texttt{check\_\allowbreak{}availability} & -- \\
T1-04 & \texttt{holding}, party 8, \texttt{next-week} & ``Forget this half-finished booking, we never completed it.'' & \texttt{discard\_\allowbreak{}draft}, \texttt{release\_\allowbreak{}hold} & -- \\
T1-05 & \texttt{waitlisted}, party 8, \texttt{saturday} & ``Take us off the queue, we have made other plans.'' & \texttt{leave\_\allowbreak{}waitlist} & -- \\
T1-06 & \texttt{details\_\allowbreak{}entered}, party 2, \texttt{next-week} & ``We would like to order the food in advance.'' & \texttt{open\_\allowbreak{}preorder} & -- \\
T1-07 & \texttt{cancellation\_\allowbreak{}requested}, party 8, \texttt{tomorrow} & ``Start again from scratch.'' & \texttt{start\_\allowbreak{}over} & -- \\
T1-08 & \texttt{browsing}, party 8, \texttt{saturday} & ``Get us in the queue for a table on that day.'' & \texttt{join\_\allowbreak{}waitlist} & -- \\
T1-09 & \texttt{holding}, party 8, \texttt{next-week} & ``Actually, we would like a different time.'' & \texttt{change\_\allowbreak{}slot} & -- \\
T1-10 & \texttt{waitlisted}, party 8, \texttt{saturday} & ``Tell me how far up the queue we are.'' & \texttt{check\_\allowbreak{}position} & -- \\
\hline
T2-01 & \texttt{browsing} & ``Get us in the queue for a table on that day.'' & \texttt{join\_\allowbreak{}waitlist} & \texttt{GET /availability} \\
T2-02 & \texttt{holding}, party 8, \texttt{next-week} & ``Go ahead and confirm the booking.'' & \texttt{confirm\_\allowbreak{}booking}, \texttt{pay\_\allowbreak{}deposit} & \texttt{POST /details} \\
T2-03 & \texttt{details\_\allowbreak{}entered}, party 2, \texttt{next-week} & ``Add the tiramisù to what we are ordering ahead.'' & \texttt{add\_\allowbreak{}dish} & \texttt{POST /preorder} \\
T2-04 & \texttt{deposit\_\allowbreak{}pending}, party 8, \texttt{next-week} & ``Actually, we would like a different time.'' & \texttt{change\_\allowbreak{}slot} & \texttt{POST /edit-details} \\
T2-05 & \texttt{browsing} & ``Put a table on hold for us.'' & \texttt{hold\_\allowbreak{}slot} & \texttt{GET /availability} \\
T2-06 & \texttt{holding}, party 2, \texttt{next-week} & ``We would like to order the food in advance.'' & \texttt{open\_\allowbreak{}preorder} & \texttt{POST /details} \\
T2-07 & \texttt{details\_\allowbreak{}entered}, party 2, \texttt{next-week} & ``Add the pasta dish to the advance order.'' & \texttt{add\_\allowbreak{}dish} & \texttt{POST /preorder} \\
T2-08 & \texttt{holding}, party 2, \texttt{next-week}, pre-order open & ``Take that dish off the advance order.'' & \texttt{remove\_\allowbreak{}dish} & \texttt{POST /details} \\
\hline
T3-01 & \texttt{holding}, party 8, \texttt{next-week} & ``Add the tiramisù to what we are ordering ahead.'' & \texttt{add\_\allowbreak{}dish} & -- \\
T3-02 & \texttt{confirmed}, party 8, \texttt{next-week} & ``Forget this half-finished booking, we never completed it.'' & \texttt{discard\_\allowbreak{}draft}, \texttt{release\_\allowbreak{}hold} & -- \\
T3-03 & \texttt{conflict}, party 8, \texttt{next-week} & ``Go ahead and confirm the booking.'' & \texttt{confirm\_\allowbreak{}booking}, \texttt{pay\_\allowbreak{}deposit} & -- \\
T3-04 & \texttt{cancellation\_\allowbreak{}requested}, party 8, \texttt{tomorrow} & ``We cannot make it — cancel the booking.'' & \texttt{cancel\_\allowbreak{}booking}, \texttt{request\_\allowbreak{}cancellation}, \texttt{discard\_\allowbreak{}draft}, \texttt{release\_\allowbreak{}hold} & -- \\
T3-05 & \texttt{cancelled}, party 8, \texttt{next-week} & ``What is happening with the cancellation we asked for?'' & \texttt{view\_\allowbreak{}cancellation} & -- \\
T3-06 & \texttt{holding}, party 8, \texttt{next-week} & ``Add the pasta dish to the advance order.'' & \texttt{add\_\allowbreak{}dish} & -- \\
\end{longtable}
}

\clearpage

\section{What Each Condition Sees}\label{app:surfaces}

This appendix shows the first model request of task T1-01 under each condition.
In the start state a party of eight has entered the guest details and must pay a deposit before the booking is confirmed (G1).
The goal is ``Go ahead and confirm the booking.''
Three operations are valid: \texttt{pay\_deposit}, \texttt{edit\_details} and \texttt{discard\_draft}.

\subsection{The page}\label{app:surfaces:page}

Every condition except P receives the same \texttt{CURRENT PAGE} block, the visible text of the page:
\begin{Verbatim}
CURRENT PAGE:
Deposit required
next-week at 20:00 for 8 — Ada Lovelace (+358 40 123 4567)
Parties of more than 6 require a €80 deposit before the booking can be confirmed.
\end{Verbatim}
P receives the raw HTML instead; this is the complete user message under P:
\begin{Verbatim}
GOAL: Go ahead and confirm the booking.

CURRENT PAGE (raw HTML):
<!DOCTYPE html>
<html lang="en">
<head>
<meta charset="utf-8">
<meta name="viewport" content="width=device-width, initial-scale=1">
<title>La Trattoria — reservations</title>
<link rel="stylesheet" href="https://cdn.jsdelivr.net/npm/@picocss/pico@2/css/pico.min.css">
<script src="https://unpkg.com/htmx.org@2.0.4"></script>
</head>
<body>
<header class="container">
<nav><ul><li><strong>La Trattoria</strong></li></ul><ul><li><small>reservations</small></li></ul></nav>
</header>
<main class="container">
<article id="booking">
<h2>Deposit required</h2>
<p><strong>next-week</strong> at <strong>20:00</strong> for <strong>8</strong> — Ada Lovelace (+358 40 123 4567)</p>
<p><mark>Parties of more than 6 require a €80 deposit
before the booking can be confirmed.</mark></p>

<form method="POST" action="/deposit">
<label>Card number<input name="card" type="text" required></label>
<button type="submit">Pay €80 deposit</button>
</form>
<form method="POST" action="/edit-details">

<button type="submit">Edit details</button>
</form>
<button type="button" hx-delete="/draft" hx-target="#booking" hx-swap="innerHTML">Discard draft</button>
</article>
</main>
</body>
</html>
\end{Verbatim}

\subsection{J and J+}\label{app:surfaces:catalogue}

The list of all 21 operations under J+.
Under J each entry is the same without its second sentence, the precondition.
\begin{Verbatim}
AVAILABLE ACTIONS:
[0] List the tables free on a given date for a given party size. Only meaningful before a table is held. — GET /availability
      - date (text, required, one of: tomorrow, next-week, saturday)
      - partySize (text, required)
[1] Put a temporary hold on a time slot. Only possible if that slot is still free; a fully booked date offers the waitlist instead. — POST /hold
      - slot (text, required)
[2] Join the queue for a table on the chosen date. Only offered when no slot on the date is free. — POST /waitlist
[3] Leave the waitlist queue. Only while queued. — DELETE /waitlist
[4] Read your current position in the waitlist queue. Only while queued. — GET /waitlist
[5] Attach the guest's name and phone number to the held table. Only while a table is held. — POST /details
      - name (text, required)
      - phone (text, required)
[6] Give up the held slot and return to the slot list. Only while a table is held and not yet confirmed. — POST /change-slot
[7] Release the temporary hold on the slot. Only while a table is held. — DELETE /hold
[8] Turn the reviewed draft into a confirmed booking. Only for parties that do not require a deposit; large parties must pay a deposit instead. — POST /confirm
[9] Pay the deposit a large party requires, which also confirms the booking. Only when the party size is above the restaurant's threshold. — POST /deposit
      - card (text, required)
[10] Go back and re-enter the guest details. Only before the booking is confirmed. — POST /edit-details
[11] Throw away the unconfirmed draft booking. Only before confirmation; this is not the same action as cancelling a confirmed booking. — DELETE /draft
[12] Start a food pre-order for the booking. Only after details are entered and before the booking is confirmed. — POST /preorder
[13] Add a dish to the pre-order. Only while the pre-order is open and not yet submitted. — POST /preorder/dishes
      - dish (text, required, one of: Bruschetta, Tagliatelle al ragù, Tiramisù)
[14] Remove a dish from the pre-order. Only while the pre-order is open and not yet submitted. — DELETE /preorder/dishes
      - dish (text, required, one of: Bruschetta, Tagliatelle al ragù, Tiramisù)
      - index (text, required)
[15] Submit the pre-order, locking it. Only while the pre-order is open. — POST /preorder/submit
[16] Read the confirmation page for the booking. Only once the booking is confirmed. — GET /confirmation
[17] Cancel a confirmed booking outright. Only when seating is more than 24 hours away; closer than that, a cancellation must be requested. — DELETE /booking
[18] Ask the restaurant to cancel a booking, giving a reason. Only within 24 hours of seating; earlier than that, the booking can simply be cancelled. — POST /cancellation-request
      - reason (text, required)
[19] Read the status of a pending cancellation request. Only after a cancellation has been requested. — GET /cancellation
[20] Abandon everything and start a new session. Only from a finished session. — GET /reset
\end{Verbatim}

\subsection{X and X+}\label{app:surfaces:mcp}

Under X the \texttt{tools} parameter carries the 21 operations as tool definitions.
Each description is the J+ entry word for word; the parameters are a JSON Schema with the same domains.
One of the 21:
\begin{Verbatim}
{
  "type": "function",
  "function": {
    "name": "pay_deposit",
    "description": "Pay the deposit a large party requires, which also confirms the booking. Only when the party size is above the restaurant's threshold.",
    "parameters": {
      "type": "object",
      "properties": {
        "card": {
          "type": "string",
          "description": "The card to use."
        }
      },
      "required": [
        "card"
      ],
      "additionalProperties": false
    }
  }
}
\end{Verbatim}
Under X+ the server returns only the tools valid in the current state, here \texttt{pay\_deposit}, \texttt{edit\_details}, \texttt{discard\_draft}, with the same definitions as under X. The text of the user message is the same as under X: the goal and the page.

\subsection{H}\label{app:surfaces:h}

The list of the controls parsed from the page:
\begin{Verbatim}
AVAILABLE ACTIONS:
[0] Pay €80 deposit — POST /deposit
      - card (text, required)
[1] Edit details — POST /edit-details
[2] Discard draft — DELETE /draft
\end{Verbatim}
